% ============================================================
%  Probeklausur Template (my_dhbw Stil, klausur_*.tex)
%  Formell mit Deckblatt-Header; Lösungen als grüne tcolorbox.
%  Renderer: pdflatex (zweimal)
% ============================================================
\documentclass[11pt, a4paper]{article}

\usepackage[ngerman]{babel}
\usepackage{fontspec}
\setmainfont[
  Path=/usr/share/fonts/TTF/,
  Extension=.ttf,
  BoldFont=DejaVuSans-Bold,
  ItalicFont=DejaVuSans-Oblique,
  BoldItalicFont=DejaVuSans-BoldOblique
]{DejaVuSans}
\setsansfont[
  Path=/usr/share/fonts/TTF/,
  Extension=.ttf,
  BoldFont=DejaVuSans-Bold,
  ItalicFont=DejaVuSans-Oblique,
  BoldItalicFont=DejaVuSans-BoldOblique
]{DejaVuSans}
\setmonofont[Path=/usr/share/fonts/TTF/,Extension=.ttf]{DejaVuSansMono}
\usepackage[top=2cm, bottom=2cm, left=2.4cm, right=2.4cm, footskip=1cm]{geometry}
\usepackage{amsmath, amssymb}
\usepackage{xcolor}
\usepackage{enumitem}
\usepackage{fancyhdr}
\usepackage{lastpage}
\usepackage{tabularx}
\usepackage{booktabs}
\usepackage{microtype}
\usepackage{parskip}

\usepackage{array}
\usepackage{longtable}
\usepackage{ltablex}
\keepXColumns
\usepackage{colortbl}
\usepackage[most,breakable]{tcolorbox}
\usepackage{listings}
\usepackage[normalem]{ulem}
\usepackage{pdflscape}
\usepackage{titlesec}
\usepackage[colorlinks=true, linkcolor=accent, urlcolor=accent]{hyperref}

\renewcommand{\familydefault}{\sfdefault}

% ── Theme ─────────────────────────────────────────────────────
\definecolor{accent}{RGB}{0, 60, 113}
\definecolor{primaryblue}{RGB}{0, 60, 113}
\definecolor{loesung}{RGB}{0, 100, 0}
\definecolor{codebg}{RGB}{245, 245, 245}
\definecolor{figurebg}{RGB}{232, 241, 255}
\definecolor{tablehintbg}{RGB}{240, 255, 240}
\definecolor{solutionbg}{RGB}{240, 252, 240}
\definecolor{rulegray}{RGB}{180, 180, 180}
\definecolor{tableheadbg}{RGB}{0, 60, 113}

% ── Header/Footer ─────────────────────────────────────────────
\pagestyle{fancy}
\fancyhf{}
\renewcommand{\headrulewidth}{0.4pt}
\fancyhead[L]{\footnotesize\color{accent}Modulprüfung -- \nichttitle}
\fancyhead[R]{\footnotesize\color{accent}\nichtsubtitle}
\fancyfoot[L]{\footnotesize\color{accent}Matrikelnr.: \rule{3cm}{0.4pt}}
\fancyfoot[R]{\footnotesize\color{accent}Blatt \thepage\,/\,\pageref{LastPage}}

% ── Typografie ────────────────────────────────────────────────
\titleformat{\section}
    {\color{accent}\large\bfseries}{}{0pt}{}
    [\vspace{-4pt}\color{accent}\rule{\linewidth}{1.5pt}]
\titleformat{\subsection}
    {\normalsize\bfseries\color{accent!85!black}}{}{0pt}{}{}
\titleformat{\subsubsection}
    {\small\bfseries\color{accent!70!black}}{}{0pt}{}{}
\titlespacing*{\section}{0pt}{12pt}{4pt}
\titlespacing*{\subsection}{0pt}{8pt}{2pt}
\titlespacing*{\subsubsection}{0pt}{6pt}{1pt}

\setlength{\parindent}{0pt}
\setlength{\parskip}{6pt}
\renewcommand{\arraystretch}{1.25}
\setlist[itemize]{noitemsep, topsep=3pt, parsep=0pt, leftmargin=1.4em}
\setlist[enumerate]{noitemsep, topsep=3pt, parsep=0pt, leftmargin=1.6em}

% ── Boxen ─────────────────────────────────────────────────────
\newtcolorbox{figurebox}[1]{
  colback=figurebg, colframe=accent!50,
  title={Abbildung: #1}, fonttitle=\small\bfseries,
  breakable, before skip=6pt, after skip=6pt
}
\newtcolorbox{tablehintbox}[1]{
  colback=tablehintbg, colframe=green!50!black!50,
  title={Tabelle: #1}, fonttitle=\small\bfseries,
  breakable, before skip=6pt, after skip=6pt
}
\newtcolorbox{solutionbox}[1]{
  enhanced,
  colback=solutionbg, colframe=loesung,
  title={#1}, fonttitle=\small\bfseries,
  coltitle=white, colbacktitle=loesung,
  breakable, before skip=8pt, after skip=8pt
}

\lstset{
  basicstyle=\ttfamily\small,
  backgroundcolor=\color{codebg},
  breaklines=true,
  frame=single, framerule=0pt,
  xleftmargin=4pt, xrightmargin=4pt,
  aboveskip=6pt, belowskip=6pt,
  keepspaces=true
}

\providecommand{\nichttitle}{Probeklausur}
\providecommand{\nichtsubtitle}{}

\renewcommand{\nichttitle}{Relationen, Algebra, Diskrete Mathematik}
\renewcommand{\nichtsubtitle}{Probeklausur 2}


\begin{document}

% Deckblatt
\thispagestyle{empty}
\begin{center}
\vspace*{1cm}
{\Huge\bfseries\color{accent}Probeklausur}\\[8pt]
{\Large\color{accent}\nichttitle}\\[4pt]
\ifx\nichtsubtitle\empty\else
  {\large\color{accent!80!black}\nichtsubtitle}\\[6pt]
\fi
\color{accent}\rule{0.5\linewidth}{1pt}\color{black}
\end{center}

\vspace{1cm}
\begin{tabularx}{\textwidth}{@{}l X@{}}
\textbf{Studiengang:}      & \rule{6cm}{0.4pt} \\[6pt]
\textbf{Matrikelnummer:}    & \rule{6cm}{0.4pt} \\[6pt]
\textbf{Name, Vorname:}    & \rule{6cm}{0.4pt} \\[6pt]
\textbf{Datum:}             & \rule{6cm}{0.4pt} \\[6pt]
\textbf{Bearbeitungszeit:}  & 60 Minuten \\[6pt]
\textbf{Hilfsmittel:}       & Taschenrechner (nicht programmierbar) \\
\end{tabularx}

\vspace{2cm}
\begin{center}
{\large\itshape Viel Erfolg!}
\end{center}

\newpage

% ── BODY ──────────────────────────────────────────────────────

\section{Probeklausur 2 — Relationen, Algebra, Diskrete Mathematik}

\textbf{Bearbeitungszeit}: 60 Minuten | \textbf{Hilfsmittel}: keine | \textbf{Punkte}: 100



\noindent\textcolor{rulegray}{\rule{\linewidth}{0.4pt}}


\paragraph{Aufgabe 1 — Eigenschaften von Relationen \textit{(18 Punkte)}}\mbox{}\\[2pt]

Auf der Menge M = \{1, 2, 3, 4\} sei die Relation

R = \{(1,1), (2,2), (3,3), (4,4), (1,2), (2,1), (3,4), (4,3), (1,3), (3,1)\}

gegeben.


a) Nennen Sie die formale Bedingung für Transitivität einer Relation R ⊆ M². \textit{(3 P)}


b) Untersuchen Sie R systematisch auf die Eigenschaften reflexiv, symmetrisch, antisymmetrisch und transitiv. Begründen Sie jede Aussage durch Verweis auf die Definition oder ein konkretes Gegenbeispiel-Tupel. \textit{(10 P)}


c) Ergänzen Sie R um die minimal nötigen Tupel, sodass R zu einer Äquivalenzrelation auf M wird. Geben Sie anschließend die entstehenden Äquivalenzklassen an. \textit{(5 P)}



\noindent\textcolor{rulegray}{\rule{\linewidth}{0.4pt}}


\paragraph{Aufgabe 2 — Ordnungsrelationen analysieren \textit{(14 Punkte)}}\mbox{}\\[2pt]

Auf der Menge S aller nichtleeren Teilmengen von \{a, b, c\} sei die Relation

A ⊑ B :⇔ A ⊆ B

definiert.


a) Begründen Sie kurz, warum ⊑ reflexiv und antisymmetrisch ist. \textit{(4 P)}


b) Klassifizieren Sie ⊑ gemäß der Tabelle Quasi-/Halb-/Vollordnung (reflexive Variante). Begründen Sie insbesondere, ob ⊑ linear ist; geben Sie ggf. zwei unvergleichbare Elemente an. \textit{(6 P)}


c) Auf der Menge ℕ sei R durch xRy :⇔ x ≥ y definiert. Ein Mitstudent behauptet, R sei eine reflexive Halbordnung. Widerlegen Sie diese Behauptung mit der passenden Eigenschaft und einem Gegenbeispiel. \textit{(4 P)}



\noindent\textcolor{rulegray}{\rule{\linewidth}{0.4pt}}


\paragraph{Aufgabe 3 — Abbildungen und ihre Eigenschaften \textit{(16 Punkte)}}\mbox{}\\[2pt]

Gegeben seien die folgenden vier Abbildungen:


\begin{itemize}
  \item f₁: ℝ → ℝ, f₁(x) = 2x − 5
  \item f₂: ℝ → ℝ⁺₀, f₂(x) = x²
  \item f₃: ℤ → ℤ, f₃(x) = ⌊x/2⌋
  \item f₄: ℝ⁺ → ℝ⁺, f₄(x) = x²
\end{itemize}

a) Erstellen Sie eine Tabelle und tragen Sie für jede Abbildung ein, ob sie linksvollständig, rechtseindeutig, injektiv, surjektiv und bijektiv ist (✓/✗). \textit{(10 P)}


b) Begründen Sie ausführlich für f₂, warum sie nicht injektiv, aber surjektiv ist. \textit{(3 P)}


c) Eine Mitstudentin sagt: „f(x) = x² ist nie bijektiv." Widerlegen Sie diese Aussage mit Bezug auf f₄ und erklären Sie, wovon die Bewertung „bijektiv" inhaltlich abhängt. \textit{(3 P)}



\noindent\textcolor{rulegray}{\rule{\linewidth}{0.4pt}}


\paragraph{Aufgabe 4 — Datenbankentwurf bewerten und korrigieren \textit{(20 Punkte)}}\mbox{}\\[2pt]

Ein Sportverein speichert seine Mitglieder in folgender Tabelle \texttt{MITGLIEDER}:



\begin{landscape}

{\footnotesize
\begin{tabularx}{\linewidth}{XXXXXXXX}
\hline
\rowcolor{tableheadbg}
{\color{white}\textbf{MNr}} & {\color{white}\textbf{Name}} & {\color{white}\textbf{Geburtsjahr}} & {\color{white}\textbf{Fußball}} & {\color{white}\textbf{Tennis}} & {\color{white}\textbf{Schwimmen}} & {\color{white}\textbf{TrainerName}} & {\color{white}\textbf{TrainerTel}} \\[2pt]
\hline
\endhead
\hline
\endfoot
1 & Anna & 2003 & 1 & 0 & 1 & Müller & 0151-… \\
2 & Ben & 2002 & 1 & 0 & 0 & Müller & 0151-… \\
3 & Clara & 2004 & 0 & 1 & 1 & Schmidt & 0162-… \\
\end{tabularx}
}
\end{landscape}


a) Identifizieren Sie zwei verschiedene Probleme dieses Entwurfs (eines bezogen auf die Sportart-Spalten, eines bezogen auf die Trainer-Daten) und benennen Sie jeweils die Anomalie-Art. \textit{(6 P)}


b) Entwerfen Sie ein normalisiertes Schema mit drei Tabellen. Geben Sie für jede Tabelle den Tabellennamen, die Attribute, den Primärschlüssel (unterstrichen oder mit „PK") und alle Fremdschlüssel an. \textit{(10 P)}


c) Begründen Sie, warum (Name, Geburtsjahr) trotz scheinbarer Eindeutigkeit ein schlechter Primärschlüsselkandidat in \texttt{MITGLIEDER} ist. Beziehen Sie sich auf die Definition „eindeutig + minimal". \textit{(4 P)}



\noindent\textcolor{rulegray}{\rule{\linewidth}{0.4pt}}


\paragraph{Aufgabe 5 — Relationenalgebra anwenden \textit{(20 Punkte)}}\mbox{}\\[2pt]

Gegeben seien die Relationen:


\texttt{STUDENT(MatrNr, Name, KursID)}


{\small
\begin{tabularx}{\linewidth}{XXX}
\hline
\rowcolor{tableheadbg}
{\color{white}\textbf{MatrNr}} & {\color{white}\textbf{Name}} & {\color{white}\textbf{KursID}} \\[2pt]
\hline
\endhead
\hline
\endfoot
101 & Alex & TINF24 \\
102 & Bea & WINF24 \\
103 & Cem & TINF24 \\
104 & Dana & WINF23 \\
\end{tabularx}
}


\texttt{KURS(KursID, Standort)}


{\small
\begin{tabularx}{\linewidth}{XX}
\hline
\rowcolor{tableheadbg}
{\color{white}\textbf{KursID}} & {\color{white}\textbf{Standort}} \\[2pt]
\hline
\endhead
\hline
\endfoot
TINF24 & Stuttgart \\
WINF24 & Mannheim \\
WINF23 & Mannheim \\
\end{tabularx}
}


a) Berechnen Sie das Ergebnis von σ\_\{Standort='Mannheim'\}(KURS) und geben Sie das Ergebnis als Tabelle an. \textit{(3 P)}


b) Berechnen Sie das Ergebnis von Π\_\{Name\}(σ\_\{KursID='TINF24'\}(STUDENT)). \textit{(4 P)}


c) Berechnen Sie STUDENT ⋈ KURS (Natural Join) und geben Sie das Ergebnis als Tabelle an. \textit{(6 P)}


d) Formulieren Sie die folgende Anfrage in Relationenalgebra: „Namen aller Studierenden, deren Kurs in Mannheim stattfindet." Geben Sie zusätzlich das äquivalente SQL-Statement mit \texttt{inner join} an. \textit{(7 P)}



\noindent\textcolor{rulegray}{\rule{\linewidth}{0.4pt}}


\paragraph{Aufgabe 6 — Mengenoperationen und Schema-Kompatibilität \textit{(12 Punkte)}}\mbox{}\\[2pt]

Zwei Tabellen sollen vereinigt werden:


\texttt{A(Nr, Bezeichnung)} mit Tupeln \{(1, 'rot'), (2, 'blau')\}

\texttt{B(ID, Farbe)} mit Tupeln \{(2, 'blau'), (3, 'grün')\}


a) Ein Studierender berechnet einfach A ∪ B = \{(1,'rot'), (2,'blau'), (3,'grün')\}. Erklären Sie, warum diese Operation in der Relationenalgebra so nicht zulässig ist. \textit{(4 P)}


b) Geben Sie zwei mögliche Vorgehensweisen an, mit denen die Vereinigung dennoch korrekt durchgeführt werden kann. \textit{(4 P)}


c) Berechnen Sie das Kreuzprodukt A × B und geben Sie an, wie viele Tupel das Ergebnis enthält. Begründen Sie, warum hier — anders als bei ∪ — keine Schema-Kompatibilität nötig ist. \textit{(4 P)}



\noindent\textcolor{rulegray}{\rule{\linewidth}{0.4pt}}



\section{Musterlösung Klausur 2}

\paragraph{Aufgabe 1 \textit{(18 Punkte)}}\mbox{}\\[2pt]

\textbf{a)} \textit{(3 P)} Transitivität: ∀x,y,z ∈ M: (xRy ∧ yRz) → xRz.

\textit{Bewertung:} 3 P bei korrekter Formel mit Quantor und Implikation; 1 P Abzug bei fehlendem Quantor.


\textbf{b)} \textit{(10 P)}

\begin{itemize}
  \item \textbf{Reflexiv ✓}: (1,1), (2,2), (3,3), (4,4) ∈ R, also xRx für alle x ∈ M. \textit{(2 P)}
  \item \textbf{Symmetrisch ✓}: Zu prüfen sind alle nicht-diagonalen Paare: (1,2)↔(2,1) ✓, (3,4)↔(4,3) ✓, (1,3)↔(3,1) ✓. \textit{(2 P)}
  \item \textbf{Antisymmetrisch ✗}: Gegenbeispiel (1,2) ∈ R und (2,1) ∈ R, aber 1 ≠ 2. \textit{(2 P)}
  \item \textbf{Transitiv ✗}: (2,1) ∈ R und (1,3) ∈ R, aber (2,3) ∉ R. (Alternativ: (4,3) und (3,1), aber (4,1) ∉ R.) \textit{(4 P)}
\end{itemize}

\textit{Bewertung:} je 2 P pro Eigenschaft inkl. Begründung; bei „transitiv ✗" 4 P, weil Gegenbeispiel anspruchsvoller.


\textbf{c)} \textit{(5 P)} Für Äquivalenz fehlt Transitivität. Aus den Symmetrie-/Verbindungspaaren folgt: 1 ist mit 2 und 3 verbunden, 3 mit 4. Damit müssen 1, 2, 3, 4 alle in einer Klasse liegen. Zu ergänzen: (2,3), (3,2), (1,4), (4,1), (2,4), (4,2). Äquivalenzklassen: nur \{1, 2, 3, 4\}. \textit{(3 P Tupel + 2 P Klasse)}



\noindent\textcolor{rulegray}{\rule{\linewidth}{0.4pt}}


\paragraph{Aufgabe 2 \textit{(14 Punkte)}}\mbox{}\\[2pt]

\textbf{a)} \textit{(4 P)}

\begin{itemize}
  \item Reflexiv: A ⊆ A gilt für jede Menge A. \textit{(2 P)}
  \item Antisymmetrisch: A ⊆ B ∧ B ⊆ A → A = B (Mengengleichheit per Definition). \textit{(2 P)}
\end{itemize}

\textbf{b)} \textit{(6 P)} ⊑ ist reflexiv, antisymmetrisch und transitiv (A ⊆ B ⊆ C → A ⊆ C). Damit ist ⊑ eine \textbf{reflexive Halbordnung}. ⊑ ist \textbf{nicht linear}: \{a\} und \{b\} sind unvergleichbar, da weder \{a\} ⊆ \{b\} noch \{b\} ⊆ \{a\}. Daher \textbf{keine Vollordnung}. \textit{(3 P Klassifikation + 3 P Begründung Linearität mit Beispiel)}


\textbf{c)} \textit{(4 P)} R = ≥ ist reflexiv (x ≥ x ✓) und transitiv ✓, aber \textbf{nicht antisymmetrisch}, sondern wäre es erst mit Antisymmetrie. Doch x ≥ y ∧ y ≥ x → x = y gilt sogar — d. h. ≥ IST antisymmetrisch. Korrekte Widerlegung: ≥ ist sogar \textbf{linear} (Vollordnung), nicht nur Halbordnung — die Aussage „nur Halbordnung" ist also zu schwach. Alternativ-akzeptiert: Hinweis aus Zettel: „größer oder gleich groß" auf Mengen-Niveau ist keine reflex. Halbordnung, weil dort Antisymmetrie fehlt; auf ℕ mit ≥ liegt jedoch sogar Vollordnung vor. \textit{(2 P Eigenschaft + 2 P Beispiel/Begründung)}


\textit{Bewertungshinweis:} Volle Punktzahl auch, wenn Studi argumentiert „R ist sogar Vollordnung, weil zusätzlich linear: ∀x,y ∈ ℕ: x ≥ y ∨ y ≥ x".



\noindent\textcolor{rulegray}{\rule{\linewidth}{0.4pt}}


\paragraph{Aufgabe 3 \textit{(16 Punkte)}}\mbox{}\\[2pt]

\textbf{a)} \textit{(10 P, je 0,5 P pro Zelle)}



{\small
\begin{tabularx}{\linewidth}{XXXXXX}
\hline
\rowcolor{tableheadbg}
{\color{white}\textbf{f}} & {\color{white}\textbf{linksvollst.}} & {\color{white}\textbf{rechtseind.}} & {\color{white}\textbf{injektiv}} & {\color{white}\textbf{surjektiv}} & {\color{white}\textbf{bijektiv}} \\[2pt]
\hline
\endhead
\hline
\endfoot
f₁ & ✓ & ✓ & ✓ & ✓ & ✓ \\
f₂ & ✓ & ✓ & ✗ & ✓ & ✗ \\
f₃ & ✓ & ✓ & ✗ & ✓ & ✗ \\
f₄ & ✓ & ✓ & ✓ & ✓ & ✓ \\
\end{tabularx}
}


\textit{Begründungs-Hinweise (intern):} f₁ linear mit Steigung ≠ 0 → bijektiv. f₂ nicht injektiv: f₂(2)=f₂(−2)=4; surjektiv auf ℝ⁺₀, da ∀y≥0 ∃x=√y. f₃: f₃(2)=f₃(3)=1, nicht injektiv; surjektiv, da ∀y∈ℤ: f₃(2y)=y. f₄ auf ℝ⁺: streng monoton wachsend → injektiv; ∀y>0: x=√y > 0, also surjektiv.


\textbf{b)} \textit{(3 P)} Nicht injektiv: f₂(−1)=f₂(1)=1, aber −1 ≠ 1, verletzt also Definition (f(x₁)=f(x₂) → x₁=x₂). Surjektiv auf ℝ⁺₀: zu jedem y ≥ 0 ist x = √y ∈ ℝ Urbild mit f₂(x) = y. \textit{(je 1,5 P)}


\textbf{c)} \textit{(3 P)} Widerlegung: f₄ ist bijektiv (s. Tabelle), also gibt es eine bijektive x²-Abbildung. \textbf{Bewertung hängt von D und W ab} (siehe Zettel: „Bewertung hängt von gewählten D, W ab"); durch Einschränkung von ℝ auf ℝ⁺ wird Injektivität erreicht, durch passende Wertemenge die Surjektivität. \textit{(2 P Widerlegung + 1 P Bezug auf D/W)}



\noindent\textcolor{rulegray}{\rule{\linewidth}{0.4pt}}


\paragraph{Aufgabe 4 \textit{(20 Punkte)}}\mbox{}\\[2pt]

\textbf{a)} \textit{(6 P)}

\begin{itemize}
  \item \textbf{Sportart-Spalten} (Fußball, Tennis, Schwimmen): Pro Sportart eine eigene Spalte → viele 0-Einträge (Speicherverschwendung); neue Sportart erfordert Schema-Änderung (Spalte hinzufügen). Klassische \textbf{Insert-Anomalie / mangelnde Erweiterbarkeit}. \textit{(3 P)}
  \item \textbf{Trainer-Daten} (TrainerName, TrainerTel): Mehrfach abgespeichert (Müller bei mehreren Mitgliedern). Bei Telefonnummern-Änderung müssen alle Zeilen geändert werden → \textbf{Update-Anomalie / Redundanz}. \textit{(3 P)}
\end{itemize}

\textbf{b)} \textit{(10 P)} Normalisiertes Schema:


\begin{lstlisting}
MITGLIED(MNr [PK], Name, Geburtsjahr, TrainerID [FK → TRAINER.TrainerID])
SPORTART(SportID [PK], Sportname)
MITGLIED_SPORT(MNr [PK,FK → MITGLIED], SportID [PK,FK → SPORTART])
TRAINER(TrainerID [PK], Name, Tel)
\end{lstlisting}


\textit{Anmerkung:} 4 Tabellen sind besser als 3 — volle Punktzahl auch bei 3-Tabellen-Lösung, falls Trainer in MITGLIED-Tabelle integriert (dann 1 P Abzug wegen verbleibender Redundanz). Mindestanforderung für volle Punktzahl: separate \texttt{MITGLIED\_SPORT}-Verknüpfungstabelle (n:m), korrekte PKs/FKs. \textit{(4 P Strukturzerlegung, 3 P PKs, 3 P FKs)}


\textbf{c)} \textit{(4 P)} Eindeutigkeit: Es kann zwei Personen mit gleichem Namen und Geburtsjahr geben (Namensvettern), → Eindeutigkeit nicht garantiert. Selbst wenn zufällig eindeutig: minimal wäre nicht erfüllt, falls bereits MNr eindeutig ist (dann ist Name+Geburtsjahr nicht-minimale Erweiterung). Außerdem semantisch instabil (Name kann sich ändern, z. B. Heirat). \textit{(2 P Eindeutigkeit + 2 P Minimalität/Stabilität)}



\noindent\textcolor{rulegray}{\rule{\linewidth}{0.4pt}}


\paragraph{Aufgabe 5 \textit{(20 Punkte)}}\mbox{}\\[2pt]

\textbf{a)} \textit{(3 P)}



{\small
\begin{tabularx}{\linewidth}{XX}
\hline
\rowcolor{tableheadbg}
{\color{white}\textbf{KursID}} & {\color{white}\textbf{Standort}} \\[2pt]
\hline
\endhead
\hline
\endfoot
WINF24 & Mannheim \\
WINF23 & Mannheim \\
\end{tabularx}
}


\textbf{b)} \textit{(4 P)} σ\_\{KursID='TINF24'\}(STUDENT) liefert Tupel (101, Alex, TINF24) und (103, Cem, TINF24). Π\_\{Name\} darauf: \textbf{\{Alex, Cem\}}. \textit{(2 P Selektion + 2 P Projektion; Duplikat-Eliminierung erwähnen)}


\textbf{c)} \textit{(6 P)} Natural Join über KursID:



{\small
\begin{tabularx}{\linewidth}{XXXX}
\hline
\rowcolor{tableheadbg}
{\color{white}\textbf{MatrNr}} & {\color{white}\textbf{Name}} & {\color{white}\textbf{KursID}} & {\color{white}\textbf{Standort}} \\[2pt]
\hline
\endhead
\hline
\endfoot
101 & Alex & TINF24 & Stuttgart \\
102 & Bea & WINF24 & Mannheim \\
103 & Cem & TINF24 & Stuttgart \\
104 & Dana & WINF23 & Mannheim \\
\end{tabularx}
}


\textit{(je 1 P pro korrektes Tupel + 2 P richtiges Schema ohne KursID-Duplikat)}


\textbf{d)} \textit{(7 P)}

\begin{itemize}
  \item Relationenalgebra: Π\_\{Name\}(σ\_\{Standort='Mannheim'\}(STUDENT ⋈ KURS))
\end{itemize}
Alternativ: Π\_\{Name\}(STUDENT ⋈ σ\_\{Standort='Mannheim'\}(KURS)) — semantisch äquivalent, oft effizienter (Selektion vor Join). \textit{(4 P)}

\begin{itemize}
  \item SQL:
\end{itemize}
\begin{lstlisting}[language=sql]
  select distinct s.Name
  from STUDENT s inner join KURS k on s.KursID = k.KursID
  where k.Standort = 'Mannheim';
\end{lstlisting}

\textit{(3 P; je 1 P für distinct/select, join-Bedingung, where-Klausel)}


Ergebnis: \{Bea, Dana\}.



\noindent\textcolor{rulegray}{\rule{\linewidth}{0.4pt}}


\paragraph{Aufgabe 6 \textit{(12 Punkte)}}\mbox{}\\[2pt]

\textbf{a)} \textit{(4 P)} ∪ verlangt \textbf{Schema-Kompatibilität}: gleiche Stelligkeit UND gleiche (typgleiche) Attribute. Hier sind die Spaltennamen unterschiedlich (\texttt{Nr} vs. \texttt{ID}, \texttt{Bezeichnung} vs. \texttt{Farbe}), formal also unterschiedliche Schemata. Die Mengenoperation ist daher in der reinen Relationenalgebra nicht definiert. \textit{(2 P Begriff Schema-Kompatibilität + 2 P konkrete Begründung)}


\textbf{b)} \textit{(4 P)} Zwei mögliche Vorgehen:

\begin{enumerate}
  \item Spalten zuvor per \textbf{Umbenennung} (ρ) angleichen, z. B. ρ\_\{ID→Nr, Farbe→Bezeichnung\}(B), dann A ∪ B'.
  \item Per \textbf{Projektion} beide auf ein gemeinsames Schema bringen (sofern Typen passen) und anschließend vereinigen.
\end{enumerate}

Außerdem zu beachten: ∪ liefert Ergebnis ohne Duplikate; (2,'blau') erscheint nur einmal. \textit{(je 2 P pro Variante; 1 Bonuspunkt für Duplikat-Hinweis möglich, max. 4 P)}


\textbf{c)} \textit{(4 P)} A × B = jedes Tupel aus A mit jedem Tupel aus B kombiniert. |A| = 2, |B| = 2 → \textbf{4 Tupel}. Schema des Ergebnisses: (Nr, Bezeichnung, ID, Farbe).



{\small
\begin{tabularx}{\linewidth}{XXXX}
\hline
\rowcolor{tableheadbg}
{\color{white}\textbf{Nr}} & {\color{white}\textbf{Bezeichnung}} & {\color{white}\textbf{ID}} & {\color{white}\textbf{Farbe}} \\[2pt]
\hline
\endhead
\hline
\endfoot
1 & rot & 2 & blau \\
1 & rot & 3 & grün \\
2 & blau & 2 & blau \\
2 & blau & 3 & grün \\
\end{tabularx}
}


Begründung Schema-Kompatibilität: Beim Kreuzprodukt werden die Schemata \textbf{nebeneinandergesetzt}, nicht identifiziert; Spalten existieren parallel. Daher ist keine Übereinstimmung der Attribute nötig — anders als bei ∪/∩/∖, die elementweise auf identischen Schemata operieren. \textit{(2 P Tupelzahl + Tabelle, 2 P Begründung)}





% ──────────────────────────────────────────────────────────────

\end{document}
